\documentclass[10pt,a4paper]{amsart}
\usepackage[margin=19mm]{geometry}
\usepackage{amsmath,amssymb,mathtools}
\usepackage[scr=rsfs,cal=euler]{mathalfa}
\usepackage{xfrac}
\usepackage[unicode,hidelinks,bookmarks=false]{hyperref}
\newcommand{\ie}{i.e.\ }
\newcommand{\cf}{cf.\ }
\newcommand{\resp}{respectively\ }
\newcommand{\Subsection}{\subsection}
\providecommand{\nobreakdash}{\nobreak\hbox{-}}
\setlength{\emergencystretch}{2em}
\multlinegap0pt
\usepackage{appendix}
\newtheorem{theorem}{Theorem}[section]
\newtheorem{corollary}{Corollary}[section]
\newtheorem{proposition}{Proposition}[section]
\newtheorem{definition}{Definition}[section]
\newtheorem{lemma}{Lemma}[section]
\newtheorem{remark}{Remark}[section]
\newtheorem{example}{Example}[section]
\newtheorem{conjecture}{Conjecture}[section]
\begin{document}
\title[The Gaussian Gabor system at the critical density is a weigh]{The Gaussian Gabor system at the critical density is a weighted lower semi frame}
\author{Elias Zikkos; Abu Dhabi, United Arab Emirates}
\date{}
\maketitle
\noindent\emph{Source and licence.} The Gaussian Gabor system at the critical density is a weighted lower semi frame. Version arXiv:2606.00764v1 (CC BY 4.0). This material is distributed under the Creative Commons Attribution 4.0 International licence, http://creativecommons.org/licenses/by/4.0/, as recorded in the arXiv OAI licence metadata for this exact version and shown on the abstract page https://arxiv.org/abs/2606.00764v1. Full rights notice: licenses/arxiv-2606.00764-CC-BY-4.0.txt.\par\smallskip
\noindent\textit{Editorial scope.} Counted unit: the original \texttt{corollary} environment \texttt{GaborWLF} at source lines 91--93 together with its complete proof at lines 95--98; the delivered document keeps source lines 52--99 so that every referenced label and every citation resolves inside it. Native editable source is the paper's own concatenated TeX; the AMS wrapper is editorial formatting only. No publisher class, upstream script or unrelated artwork is redistributed.
\section{Introduction and Results}
\setcounter{equation}{0}

Let $G(\varphi, \mathbb{Z}^2)$ be the Gaussian Gabor system $G(\varphi, \mathbb{Z}^2)$ at the critical density, that is
\[
G(\varphi, \mathbb{Z}^2)=\{e^{2\pi int}\cdot e^{-\pi(t-m)^2}:\, (m,n)\in \mathbb{Z}^2\}.
\]
It is well known that $G(\varphi, \mathbb{Z}^2)$ has excess equal to one in $L^2 (\mathbb{R})$, that is,
it is complete and remains complete if any one element is removed but not if two.
It is also a Bessel sequence (upper frame) but not a frame. If any element is removed, then the remaining complete system is not a Schauder basis.
\begin{remark}
These classical facts can be found in Christensen \cite[Corollary 11.4.3 and Theorem 11.6.1]{Christensen}  
and Heil \cite[Example 11.34]{Heil}.
\end{remark}

It was recently proved in \cite[Section 6]{Balazs2026JFAA} that $G(\varphi, \mathbb{Z}^2)$ is $not$ a weighted
frame for $L^2 (\mathbb{R})$, relying on the fact, established in \cite[Corollaries 2.6 and 4.3]{Heil2025JFAA},\emph{} 
that $G(\varphi, \mathbb{Z}^2)$ does not have a reproducing partner.
The question whether it could be a $weighted\,\, lower\,\, semi\,\, frame$
for $L^2 (\mathbb{R})$ was left open. We answer this in the affirmative.
In fact, if we delete any atom from the system, the following is true.

\begin{theorem}\label{weightedlower}
For every fixed pair $(a,b)\in\mathbb{Z}^2$, consider the Gabor system
\[
\{e^{2\pi int}\cdot e^{-\pi(t-m)^2}:\, (m,n)\in \mathbb{Z}^2\setminus\{(a,b)\}\}
\]
which is complete and minimal in $L^2 (\mathbb{R})$. Let $\{\gamma_{m,n,a,b}:\, (m,n)\in \mathbb{Z}^2\setminus\{(a,b)\}\}$ be the unique biorthogonal family to this Gabor system in $L^2(\mathbb{R})$. Then
there exist positive constants $c_{m,n}$, so that the weighted Gabor system
\[
\{c_{m,n}\cdot  e^{2\pi int}\cdot e^{-\pi(t-m)^2}:\, (m,n)\in \mathbb{Z}^2\setminus\{(a,b)\}\}
\]
is a complete Riesz-Fischer sequence in $L^2 (\mathbb{R})$ and a lower semi frame for $L^2(\mathbb{R})$.
In particular this holds if we choose the constants $c_{m,n}$ so that
\[
\sum_{(m,n)\in \mathbb{Z}^2\setminus\{(a,b)\}}\frac{||\gamma_{m,n,a,b}||_{L^2(\mathbb{R})}^2}{c_{m,n}^2}<\infty.
\]
\end{theorem}


\begin{corollary}\label{GaborWLF}
The Gaussian Gabor system $G(\varphi, \mathbb{Z}^2)$ is a weighted lower semi frame for $L^2 (\mathbb{R})$.
\end{corollary}

\begin{proof}
Since $G(\varphi, \mathbb{Z}^2)$ has excess equal to one, deleting any atom yields a complete and minimal system in $L^2 (\mathbb{R})$.
By Theorem \ref{weightedlower}, suitable weights make this subsystem a complete Riesz–Fischer sequence and lower semi frame. Hence the full Gaussian Gabor system is a weighted lower semi frame.
\end{proof}


\begin{thebibliography}{99}
\bibitem[Balazs2026JFAA]{Balazs2026JFAA} P. Balazs, R. Corso, D. Stoeva, \textit{Weighted frames, weighted lower semi frames and unconditionally convergent multipliers}, J. Fourier Anal. Appl. $\bf 32$ (2026), article 24.
\bibitem[Christensen]{Christensen} O. Christensen, \textit{An Introduction to Frames and Riesz Bases}, Second Edition, Applied and Numerical Harmonic Analysis, Birkh\"{a}user, Boston, (2016).
\bibitem[Heil]{Heil} C. Heil, \textit{A Basis Theory Primer}, Expanded Edition, Applied and Numerical Harmonic Analysis, Birkh\"{a}user/Springer, New York, (2011).
\bibitem[Heil2025JFAA]{Heil2025JFAA} L. Hart, C. Heil, I. Katz, M. Northington V, \textit{Overcomplete reproducing pairs}, J. Fourier Anal. Appl. $\bf 31$ (2025), article 58
\end{thebibliography}
\end{document}
